\par\Needspace{14\baselineskip}
\SourceLeaf{139}
\begin{center}\textit{ERSTER ABSCHNITT.}\\[4pt]
\textit{Beziehungen zwischen benachbarten Funktionen:}\end{center}
\SixthArticle{7}
Wir nennen eine \textit{Funktion} zu $F(\alpha,\beta,\gamma,x)$ \textit{benachbart}, wenn sie daraus entsteht, indem das erste, zweite oder dritte Element um eins vergrößert oder verkleinert wird, während die drei übrigen Elemente unverändert bleiben. Die ursprüngliche Funktion $F(\alpha,\beta,\gamma,x)$ liefert somit sechs benachbarte Funktionen; zwischen je zweien von ihnen und der ursprünglichen Funktion selbst besteht eine sehr einfache lineare Gleichung. Diese fünfzehn Gleichungen stellen wir hier zusammen, wobei wir der Kürze halber das vierte Element, das stets als $=x$ mitzudenken ist, weglassen und die ursprüngliche Funktion einfach mit $F$ bezeichnen.
\[
\begin{array}{rcl}
{[1]}&0=&(\gamma-2\alpha-(\beta-\alpha)x)F+\alpha(1-x)F(\alpha+1,\beta,\gamma)\\
&&{}-(\gamma-\alpha)F(\alpha-1,\beta,\gamma)
\end{array}
\]
\[
\begin{array}{rcl}
{[2]}&0=&(\beta-\alpha)F+\alpha F(\alpha+1,\beta,\gamma)-\beta F(\alpha,\beta+1,\gamma)
\end{array}
\]
\[
\begin{array}{rcl}
{[3]}&0=&(\gamma-\alpha-\beta)F+\alpha(1-x)F(\alpha+1,\beta,\gamma)-(\gamma-\beta)F(\alpha,\beta-1,\gamma)
\end{array}
\]
\[
\begin{array}{rcl}
{[4]}&0=&\gamma(\alpha-(\gamma-\beta)x)F-\alpha\gamma(1-x)F(\alpha+1,\beta,\gamma)\\
&&{}+(\gamma-\alpha)(\gamma-\beta)xF(\alpha,\beta,\gamma+1)
\end{array}
\]
\[
\begin{array}{rcl}
{[5]}&0=&(\gamma-\alpha-1)F+\alpha F(\alpha+1,\beta,\gamma)-(\gamma-1)F(\alpha,\beta,\gamma-1)
\end{array}
\]
\[
\begin{array}{rcl}
{[6]}&0=&(\gamma-\alpha-\beta)F-(\gamma-\alpha)F(\alpha-1,\beta,\gamma)+\beta(1-x)F(\alpha,\beta+1,\gamma)
\end{array}
\]
\[
\begin{array}{rcl}
{[7]}&0=&(\beta-\alpha)(1-x)F-(\gamma-\alpha)F(\alpha-1,\beta,\gamma)+(\gamma-\beta)F(\alpha,\beta-1,\gamma)
\end{array}
\]
\[
\begin{array}{rcl}
{[8]}&0=&\gamma(1-x)F-\gamma F(\alpha-1,\beta,\gamma)+(\gamma-\beta)xF(\alpha,\beta,\gamma+1)
\end{array}
\]
\[
\begin{array}{rcl}
{[9]}&0=&(\alpha-1-(\gamma-\beta-1)x)F+(\gamma-\alpha)F(\alpha-1,\beta,\gamma)\\
&&{}-(\gamma-1)(1-x)F(\alpha,\beta,\gamma-1)
\end{array}
\]
\[
\begin{array}{rcl}
{[10]}&0=&(\gamma-2\beta+(\beta-\alpha)x)F+\beta(1-x)F(\alpha,\beta+1,\gamma)\\
&&{}-(\gamma-\beta)F(\alpha,\beta-1,\gamma)
\end{array}
\]
\[
\begin{array}{rcl}
{[11]}&0=&\gamma(\beta-(\gamma-\alpha)x)F-\beta\gamma(1-x)F(\alpha,\beta+1,\gamma)\\
&&{}-(\gamma-\alpha)(\gamma-\beta)F(\alpha,\beta,\gamma+1)
\end{array}
\]
\[
\begin{array}{rcl}
{[12]}&0=&(\gamma-\beta-1)F+\beta F(\alpha,\beta+1,\gamma)-(\gamma-1)F(\alpha,\beta,\gamma-1)
\end{array}
\]
\[
\begin{array}{rcl}
{[13]}&0=&\gamma(1-x)F-\gamma F(\alpha,\beta-1,\gamma)+(\gamma-\alpha)xF(\alpha,\beta,\gamma+1)
\end{array}
\]
\[
\begin{array}{rcl}
{[14]}&0=&(\beta-1-(\gamma-\alpha-1)x)F+(\gamma-\beta)F(\alpha,\beta-1,\gamma)\\
&&{}-(\gamma-1)(1-x)F(\alpha,\beta,\gamma-1)
\end{array}
\]
\[
\begin{array}{rcl}
{[15]}&0=&\gamma(\gamma-1-(2\gamma-\alpha-\beta-1)x)F\\
&&{}+(\gamma-\alpha)(\gamma-\beta)xF(\alpha,\beta,\gamma+1)\\
&&{}-\gamma(\gamma-1)(1-x)F(\alpha,\beta,\gamma-1)
\end{array}
\]
\SixthArticle{8}
Nun zum Beweis dieser Formeln. Setzt man
\[
\frac{(\alpha+1)(\alpha+2)\mathinner{.\,.\,.}(\alpha+m-1)\beta(\beta+1)\mathinner{.\,.\,.\,.\,.}(\beta+m-2)}{1\mathbin{.}2\mathbin{.}3\mathinner{.\,.\,.\,.\,.}m\mathbin{.}\gamma(\gamma+1)\mathinner{.\,.\,.\,.\,.}(\gamma+m-1)}=M
\]
\SourceLeaf{140}
so ist der Koeffizient der Potenz $x^m$
\[
\begin{array}{lll}
\mbox{in }F&.\quad.\quad.\quad.\quad.\quad.&\alpha(\beta+m-1)M\\
\mbox{in }F(\alpha,\beta-1,\gamma)&.\quad.&\alpha(\beta-1)M\\
\mbox{in }F(\alpha+1,\beta,\gamma)&.\quad.&(\alpha+m)(\beta+m-1)M\\
\mbox{in }F(\alpha,\beta,\gamma-1)&.\quad.&\frac{\alpha(\beta+m-1)(\gamma+m-1)M}{\gamma-1}
\end{array}
\]
der Koeffizient der Potenz $x^{m-1}$ in $F(\alpha+1,\beta,\gamma)$ dagegen, oder der Koeffizient der Potenz $x^m$ in $xF(\alpha+1,\beta,\gamma)$,
\[
=m(\gamma+m-1)M
\]
Daraus folgt unmittelbar die Richtigkeit der Formeln 5 und 3. Vertauscht man $\alpha$ mit $\beta$, so entsteht aus 5 die Formel 12 und aus diesen beiden durch Elimination 2. Ebenso entsteht durch dieselbe Vertauschung aus 3 die Formel 6; die Verbindung von 6 und 12 ergibt 9, daraus entsteht durch Vertauschung 14, und die Verbindung dieser beiden ergibt 7. Schließlich erhält man aus 2 und 6 die Formel 1 und daraus durch Vertauschung 10. Die Formel 8 lässt sich ähnlich wie oben die Formeln 5 und 3 durch Betrachtung der Koeffizienten herleiten (auf dieselbe Weise ließen sich, wenn man wollte, \textit{alle} 15 Formeln gewinnen), oder eleganter aus den bereits bekannten auf folgende Weise. Ersetzt man in Formel 5 das Element $\alpha$ durch $\alpha-1$ und $\gamma$ durch $\gamma+1$, so ergibt sich
\[
0=(\gamma-\alpha+1)F(\alpha-1,\beta,\gamma+1)+(\alpha-1)F(\alpha,\beta,\gamma+1)-\gamma F(\alpha-1,\beta,\gamma)
\]
Ersetzt man dagegen in Formel 9 nur $\gamma$ durch $\gamma+1$, so ergibt sich
\[
\begin{aligned}0={}&(\alpha-1-(\gamma-\beta)x)F(\alpha,\beta,\gamma+1)\\
&{}+(\gamma-\alpha+1)F(\alpha-1,\beta,\gamma+1)-\gamma(1-x)F(\alpha,\beta,\gamma)\end{aligned}
\]
Durch Subtraktion dieser Formeln entsteht unmittelbar 8 und daraus durch Vertauschung 13. Aus 1 und 8 ergibt sich 4 und daraus durch Vertauschung 11. Schließlich wird aus 8 und 9 die Formel 15 hergeleitet.
\SixthArticle{9}
Sind $\alpha'-\alpha$, $\beta'-\beta$, $\gamma'-\gamma$ sowie $\alpha''-\alpha$, $\beta''-\beta$, $\gamma''-\gamma$ ganze Zahlen (positiv oder negativ), so kann man von der Funktion $F(\alpha,\beta,\gamma)$ zur Funktion $F(\alpha',\beta',\gamma')$ und ebenso von dieser bis zur Funktion $F(\alpha'',\beta'',\gamma'')$ durch eine Folge ähnlicher Funktionen übergehen, von denen jede ihrer Vorgängerin und Nachfolgerin benachbart ist. Dazu wird zuerst ein Element, etwa $\alpha$, fortlaufend um eins verändert, bis man von $F(\alpha,\beta,\gamma)$ zu $F(\alpha',\beta,\gamma)$ gelangt ist; dann wird das zweite Element verändert, bis\SourceLeaf{141}
$F(\alpha',\beta',\gamma)$ erreicht ist, und schließlich das dritte, bis $F(\alpha',\beta',\gamma')$ erreicht ist; und ebenso geht man von dieser bis zu $F(\alpha'',\beta'',\gamma'')$ über. Da es nach Artikel 7 lineare Gleichungen zwischen der ersten, zweiten und dritten Funktion und allgemein zwischen je drei aufeinanderfolgenden Funktionen dieser Folge gibt, sieht man leicht, dass sich daraus durch Elimination eine lineare Gleichung zwischen den Funktionen $F(\alpha,\beta,\gamma)$, $F(\alpha',\beta',\gamma')$, $F(\alpha'',\beta'',\gamma'')$ herleiten lässt. Daher lässt sich, allgemein gesprochen, aus zwei Funktionen, deren erste drei Elemente sich um ganze Zahlen unterscheiden, jede andere Funktion mit derselben Eigenschaft ableiten, sofern das vierte Element dasselbe bleibt. Hier genügt es uns jedoch, diese bemerkenswerte Wahrheit allgemein festgestellt zu haben; wir verweilen hier nicht bei Abkürzungen, durch die sich die dazu nötigen Operationen möglichst kurz gestalten lassen.
\SixthArticle{10}
Es seien zum Beispiel die Funktionen
\[
F(\alpha,\beta,\gamma),\ F(\alpha+1,\beta+1,\gamma+1),\ F(\alpha+2,\beta+2,\gamma+2)
\]
gegeben, zwischen denen eine lineare Gleichung gefunden werden soll. Verbinden wir sie durch benachbarte Funktionen auf folgende Weise:
\[
\begin{array}{rl}
F(\alpha,\beta,\gamma)&=F
\end{array}
\]
\[
\begin{array}{rl}
F(\alpha+1,\beta,\gamma)&=F'
\end{array}
\]
\[
\begin{array}{rl}
F(\alpha+1,\beta+1,\gamma)&=F''
\end{array}
\]
\[
\begin{array}{rl}
F(\alpha+1,\beta+1,\gamma+1)&=F'''
\end{array}
\]
\[
\begin{array}{rl}
F(\alpha+2,\beta+1,\gamma+1)&=F''''
\end{array}
\]
\[
\begin{array}{rl}
F(\alpha+2,\beta+2,\gamma+1)&=F'''''
\end{array}
\]
\[
\begin{array}{rl}
F(\alpha+2,\beta+2,\gamma+2)&=F''''''
\end{array}
\]
Wir haben somit fünf lineare Gleichungen (aus den Formeln 6, 13, 5 des Artikels 7):
\[
\begin{array}{rcl}
\mathrm{I.}&0=&(\gamma-\alpha-1)F-(\gamma-\alpha-1-\beta)F'-\beta(1-x)F''
\end{array}
\]
\[
\begin{array}{rcl}
\mathrm{II.}&0=&\gamma F'-\gamma(1-x)F''-(\gamma-\alpha-1)xF'''
\end{array}
\]
\[
\begin{array}{rcl}
\mathrm{III.}&0=&\gamma F''-(\gamma-\alpha-1)F'''-(\alpha+1)F''''
\end{array}
\]
\[
\begin{array}{rcl}
\mathrm{IV.}&0=&(\gamma-\alpha-1)F'''-(\gamma-\alpha-2-\beta)F''''-(\beta+1)(1-x)F'''''
\end{array}
\]
\[
\begin{array}{rcl}
\mathrm{V.}&0=&(\gamma+1)F''''-(\gamma+1)(1-x)F'''''-(\gamma-\alpha-1)xF''''''
\end{array}
\]
Aus I und II ergibt sich durch Elimination von $F'$

\SourceLeaf{142}
\[
\mathrm{VI.}\quad 0=\gamma F-\gamma(1-x)F''-(\gamma-\alpha-\beta-1)xF'''
\]
Hieraus und aus III ergibt sich durch Elimination von $F''$
\[
\mathrm{VII.}\quad 0=\gamma F-(\gamma-\alpha-1-\beta x)F'''-(\alpha+1)(1-x)F''''
\]
Ferner ergibt sich aus IV und V durch Elimination von $F'''''$
\[
\mathrm{VIII.}\quad 0=(\gamma+1)F'''-(\gamma+1)F''''+(\beta+1)xF''''''
\]
Hieraus und aus VII ergibt sich durch Elimination von $F''''$
\[
\begin{aligned}&\mathrm{IX.}\quad 0=\gamma(\gamma+1)F\\&{}-(\gamma+1)(\gamma-(\alpha+\beta+1)x)F'''\\&{}-(\alpha+1)(\beta+1)x(1-x)F''''''\end{aligned}
\]
